\section{Chapter 4: Tactics}\label{sec:14-appendix-solutions:03-chapter-4:1}

\textbf{1. \texttt{theorem and\_\allowbreak{}comm\_\allowbreak{}tac \{P Q : Prop\} (h : P ∧ Q) : Q ∧ P}}

\begin{lstlisting}
theorem and_comm_tac {P Q : Prop} (h : P ∧ Q) : Q ∧ P := by
  constructor
  · exact h.right
  · exact h.left
\end{lstlisting}

\href{https://lean-lang.org/doc/reference/latest/Tactic-Proofs/Tactic-Reference/}{\texttt{constructor}} splits the goal \texttt{Q ∧ P} into two subgoals, \texttt{Q} and \texttt{P}, one for each field of \texttt{And}. \texttt{h.right : Q} closes the first, and \texttt{h.left : P} closes the second.

\textbf{2. \texttt{theorem nat\_\allowbreak{}mul\_\allowbreak{}zero (n : Nat) : n * 0 = 0}}

\begin{lstlisting}
theorem nat_mul_zero (n : Nat) : n * 0 = 0 := by
  rfl
\end{lstlisting}

\href{https://lean-lang.org/doc/reference/latest/Tactic-Proofs/Tactic-Reference/}{\texttt{rfl}} does succeed here. \texttt{Nat.mul} is defined by recursion on its second argument, and \texttt{n * 0 = 0} is the base clause. Hence this holds by definition, with no induction required. Compare this with \texttt{0 * n = 0}, which is not a base clause and does require induction on \texttt{n}.

\textbf{3. \texttt{modus\_\allowbreak{}ponens} in tactic mode}

\begin{lstlisting}
theorem modus_ponens_tac {P Q : Prop} (hpq : P → Q) (hp : P) : Q := by
  apply hpq
  exact hp
\end{lstlisting}

\texttt{apply hpq} matches the goal \texttt{Q} against the conclusion of \texttt{hpq : P → Q}. This leaves a new goal \texttt{P}, which \texttt{exact hp} closes.
